% !TeX root = ../教材题库_全部条目_转写工作稿.tex
% Source: XJ-B1-C04-HW-P147-Y0649
% 第4章 / 4.5 函数模型及其应用 / 教材 p.147 / PDF 第154页
% OCR未校对：公式与文字均可能有误；未完成数学TeX转写；禁止直接发给学生
% 原题及相邻原解截图（核对用）：../题目原图/XJ-B1-C04-HW-P147-Y0649-01.webp
% 本文件只在转写完成后作为题干使用；OCR 条目尚不是可靠题干。
\begin{bankitem}{习题4.5 第2题}{第4章 / 4.5 函数模型及其应用 / 教材 p.147 / PDF 第154页}{XJ-B1-C04-HW-P147-Y0649}{OCR未校对：公式与文字均可能有误；未完成数学TeX转写；禁止直接发给学生}
\begin{ocrdraft}
[教材p.147 / PDF第154页]
\par
2．三个变量，，3，3随着变量工的变化情况如下表，则关于分别呈对数
\par
函数、指数函数、幂函数变化的变量依次为（
\par
1
\par
3
\par
5
\par
7
\par
9
\par
11
\par
5
\par
135
\par
625
\par
17巧
\par
3645
\par
6655
\par
5
\par
29
\par
245
\par
2189
\par
19685
\par
177149
\par
Y2
\par
5
\par
6．10
\par
6．61
\par
6．985
\par
（D)YI，，3，2
\end{ocrdraft}
\end{bankitem}
